\chapteropen{10}

\begin{cele}
\item ułożyć plan na dwanaście miesięcy, w którym każdy miesiąc ma cel, rodzaj książki i jedno ćwiczenie nawyku;
\item prowadzić rejestr postępów, który liczy minuty, strony i porzucenia, a nie tylko skończone książki;
\item rozpoznać trzy kryzysy, czyli miesiąc bez czytania, nudę i trudną książkę, i przypisać każdemu plan działania;
\item przeprowadzić grudniową rewizję, odpowiadając na pięć pytań na podstawie własnego rejestru;
\item wybrać realistyczny cel roczny z tablicy i wiedzieć, kiedy wolno go podnieść.
\end{cele}

\section{Plan na dwanaście miesięcy}

Rok jest dobrą jednostką planowania, bo mieści się w nim pełny cykl świąt, urlopów, sezonów i wahań energii. Plan na rok nie polega jednak na tym, żeby z góry wybrać każdą książkę. Chodzi o rytm. Każdy miesiąc ma jeden cel liczbowy, jeden rodzaj lektury do wyboru i jedno ćwiczenie nawyku, które wykonujesz przez cały miesiąc, zanim przejdziesz dalej.

Tabela poniżej jest punktem wyjścia, a nie obietnicą. Cele są celowo skromne: w styczniu i lutym po jednej książce, potem po dwie. Rodzaje lektur przeplatają się, bo jednostajność jest częstym powodem, dla którego plany pękają w połowie roku. \verify{jednostajność lektur i nuda rzadko są jedyną przyczyną rezygnacji, ale są jedną z najczęściej opisywanych} Ćwiczenia nawyku odwołują się do rozdziałów 3, 6 i 8, bo tam ustalone zostały okna czasowe, wyzwalacze i notowanie.

\begin{center}\small
\begin{tabular}{@{}l l p{4.0cm} p{6.6cm}@{}}
\toprule
\textsf{\textbf{Miesiąc}} & \textsf{\textbf{Cel}} & \textsf{\textbf{Rodzaj książki}} & \textsf{\textbf{Ćwiczenie nawyku}}\\
\midrule
styczeń & 1 książka & literatura faktu w dziedzinie, którą już znasz & stałe okno: ta sama pora, pięć dni w tygodniu (rozdz.~3)\\
luty & 1 książka & powieść krótka, do 250 stron & wyzwalacz: czytasz zaraz po czynności, którą wykonujesz codziennie (rozdz.~6)\\
marzec & 2 książki & eseje albo popularnonaukowa książka o jednym temacie & plan awaryjny: minimum pięciu minut także w dni bez ochoty (rozdz.~6)\\
kwiecień & 2 książki & jedna pozycja porzucona wcześniej, wracająca po zasadzie pięćdziesięciu stron (rozdz.~4) & sprawdzenie zasady: dziesięć minut, a potem decyzja\\
maj & 2 książki & biografia albo historia & jedno zdanie po każdej skończonej książce (rozdz.~8)\\
czerwiec & 2 książki & jedna książka jako audiobook w podróży (rozdz.~5) & zamknięcie sesji słuchania jednym zdaniem o tym, co zostało\\
lipiec & 2 książki & tekst z drugiego zawodu: ktoś pisze o tym, co robisz w pracy & wyjaśnij komuś jeden rozdział własnymi słowami (rozdz.~8)\\
sierpień & 2 książki & trudna książka, czytana z ołówkiem (rozdz.~7) & bloki po dziesięć minut zamiast jednej długiej sesji\\
wrzesień & 2 książki & powieść, najlepiej taka, o której nikt ci nie mówił & jedno spotkanie w miesiącu z kimś, kto czyta (rozdz.~9)\\
październik & 2 książki & książka z listy z rozdziału 4 & przeniesienie sesji na inną godzinę na jeden tydzień i obserwacja\\
listopad & 2 książki & pozycja, z którą autor się nie zgadza & marginalia: jedno pytanie na stronę, najwyżej jedno zakreślenie (rozdz.~8)\\
grudzień & 2 książki & jedna krótka książka dla przyjemności i jedna dowolna & rewizja roku (sekcja niżej), zanim ustalisz cel na kolejny rok\\
\bottomrule
\end{tabular}
\end{center}

Suma celów to 22 książki. Liczba ta nie jest przypadkowa, ale wyjaśnię ją dopiero w sekcji o celu rocznym. Na razie zauważ, że plan nie zakłada, że wszystko pójdzie zgodnie z tabelą. Zakłada, że w każdym miesiącu masz coś do sprawdzenia. Jeśli w sierpniu przeczytasz połowę tego, co zaplanowałeś, to nie jest porażka planu. To jest dane, które trafia do rejestru.

\begin{zasada}
Plan roczny ustala rytm, a nie wyniki. Cel miesiąca jest liczbą, którą chcesz osiągnąć w dobrym scenariuszu, a ćwiczenie nawyku jest czynnością, którą wykonujesz niezależnie od wyniku.
\end{zasada}

\section{Śledzenie postępów}

Bez zapisu każdy rok wygląda z pamięci lepiej albo gorzej, niż był. Dlatego prowadzisz rejestr, a szablon i sposób jego prowadzenia opisuje Dodatek C. Rejestr ma sześć pól na każdą sesję lub każdą książkę: datę, tytuł, czas sesji w minutach, liczbę przeczytanych stron, decyzję (\glqq dokończona\grqq{}, \glqq porzucona\grqq{}, \glqq odłożona\grqq{}) i jedno zdanie o tym, co zaszło.

Dlaczego minuty i strony, skoro cel brzmi \emph{książki}? Bo liczba książek ukrywa wszystko, co dzieje się między skończeniami. Dwie krótkie książki mogą oznaczać dziewięćdziesiąt minut czytania, a jedna długa, zaczęta w styczniu, może stać na półce do lipca. Minuty i strony pokazują, czy nawyk żyje, nawet kiedy żadna książka nie została skończona.

\begin{przyklad}
Tydzień, w którym czytasz cztery razy po dwadzieścia pięć minut, to 100 minut. Przy 250 słowach na minutę to 25\,000 słów. Typowa książka ma \val{typical.words} słów, więc ten tydzień daje $25\,000 / 80\,000 = 0{,}31$ książki. Po miesiącu takich tygodni masz wynik, który da się porównać z celem, a nie wrażenie, że \glqq ostatnio mało czytam\grqq{}.
\end{przyklad}

\begin{falle}[Pułapka]
Najgorsza wersja rejestru zlicza tylko skończone książki. Wtedy każde porzucenie wygląda jak porażka, a każda łatwa książka jak sukces. Czytelnik zaczyna wybierać lektury według tego, jak szybko da się je zamknąć, a nie według tego, czego potrzebuje. Porzucenie jest w rejestrze zwykłą decyzją z powodem, a nie karą. Zapisuj je i wracaj do nich w grudniu.
\end{falle}

\section{Kryzysy: miesiąc bez czytania, nuda i trudna książka}

Każdy plan roczny w pewnym momencie przestaje działać. Zamiast liczyć na to, że nie wystąpi, przygotuj się do trzech konkretnych sytuacji. Każda ma inną przyczynę i inne lekarstwo, więc nie należy ich mieszać.

\subsection{Miesiąc bez czytania}

Zdarza się, że w ciągu miesiąca nie czytasz wcale: choroba, wyjazd, sesja w pracy, żałoba. Wtedy największym zagrożeniem nie jest sam brak czytania, lecz próba nadrobienia. Czytelnik, który wraca i próbuje przeczytać zaległe trzysta stron w tydzień, zwykle rezygnuje po dwóch dniach i zostaje z poczuciem winy, które jest gorsze od przerwy.

Plan jest prosty. Pierwszego dnia po powrocie czytasz minimum: pięć minut albo jedną stronę, cokolwiek wykonalne bez wysiłku. Nie nadrabiasz. Nie przenosisz zaległego celu na następny miesiąc. W rejestrze zapisujesz przyczynę przerwy jednym słowem, bo za rok będzie to przydatniejsze niż wyrzuty sumienia.

\subsection{Nuda}

Nuda ma dwa źródła i trzeba je rozdzielić. Pierwsze to książka, która jest zwyczajnie słaba albo nie dla ciebie. Drugie to nuda w czytaniu, czyli zmęczenie formą, które mija po zmianie formatu. Pierwszy przypadek wymaga porzucenia, drugi zmiany papieru na audio albo odwrotnie.

Reguła ma charakter kontrolny: jeśli dwie kolejne sesje nie dają postępu, który umiesz zapisać, zatrzymujesz się i decydujesz. Możesz zastosować zasadę pięćdziesięciu stron z rozdziału 4, zmienić format albo odłożyć książkę. Wszystkie trzy decyzje są poprawne. Błędna jest czwarta: czytanie dalej, bo się zaczęło.

\begin{falle}[Pułapka]
Najczęstsza wersja tej pułapki to \glqq przeczytam do końca, bo już tyle przeczytałem\grqq{}. Wysiłek włożony w przeszłość nie zmienia tego, ile wysiłku kosztuje następna strona. Jeśli książka jest nudna, zapisz to w rejestrze jako decyzję, a nie jako przegraną walkę.
\end{falle}

\subsection{Trudna książka}

Trudna książka to coś innego niż nudna. Czytasz ją uważnie, zatrzymujesz się, wracasz do wcześniejszych stron, ale postęp jest wolny i nie wygląda dobrze w rejestrze. Dlatego trudną książkę planuj jako jedną, a nie dwie naraz, i najlepiej obok lżejszej pozycji, jak w sierpniu z tabeli.

Pracuj w krótkich blokach po dziesięć minut, z ołówkiem, i zapisuj przy każdym bloku jedno pytanie, na które nie znasz odpowiedzi. Nie oczekuj prędkości z wcześniejszych miesięcy. Prędkość zależy od tekstu, a nie od twojego charakteru, i ten sam czytelnik czyta różnie tę samą stronę w różnych dniach. Liczy się to, że pytania są coraz lepsze, a nie to, ile stron przybywa w tygodniu.

\begin{uwaga}
Trudna książka może być jedyną, która zmienia twoje myślenie w tym roku. Nie traktuj jej jako dowodu, że nie dajesz rady, bo to zwykle dowód, że wybrałeś tekst wymagający. Porzucenie jest dopuszczalne, ale powinno być decyzją po świadomym sprawdzeniu, a nie ucieczką po pierwszym zgrzycie.
\end{uwaga}

\section{Rewizja w grudniu: pięć pytań}

W grudniu zamykasz rok i zanim ustalisz cel na kolejny, odpowiadasz na pięć pytań. Pytania są oparte na rejestrze, a nie na pamięci, bo pamięć z listopada i z grudnia jest zwykle dobra w wyolbrzymianiu ostatnich tygodni.

\begin{enumerate}
\item Ile książek skończyłeś, ile porzuciłeś i z jakim powodem zapisanym w rejestrze?
\item Które sesje były najbardziej regularne: jaki dzień tygodnia, jaka pora, jakie miejsce?
\item Które kryzysy wystąpiły i który element planu awaryjnego faktycznie zadziałał?
\item Czy jakaś książka zmieniła twoje działanie? Jeśli tak, który pomysł i jaki pierwszy krok?
\item Co zmienisz w przyszłym roku: cel roczny, okno czasowe albo rodzaje książek, i dlaczego akurat to?
\end{enumerate}

Pytania są w tej kolejności celowo. Najpierw liczby, potem wzorzec, potem zdarzenia, a dopiero na końcu decyzja. Zmiana celu zrobiona przed przejrzeniem rejestru jest zwykle zmianą nastroju, a nie wnioskiem.

\begin{przyklad}
Czytelnik, który w rejestrze ma 20 sesji po 30 minut, z czego 14 w poniedziałki i środy po obiedzie, odpowiada na drugie pytanie jasno: okno działa, a reszta tygodnia nie. Decyzja na przyszły rok nie brzmi \glqq czytać więcej\grqq{}, tylko \glqq utrzymać poniedziałki i środy, a w piątek dodać jedną krótką sesję\grqq{}. Pytanie zamieniło ogólne postanowienie w konkretne okno.
\end{przyklad}

\section{Realistyczny cel roczny}

Cel roczny wybierasz z tablicy, a nie z marzeń. Poniższa tablica pokazuje, ile książek mieści się w roku przy danej liczbie minut dziennie i przy danej szybkości czytania. Typowa książka ma \val{typical.words} słów, czyli około \val{typical.hours} godziny przy 250 słowach na minutę.

\generated{books-per-year}

Przy trzydziestu minutach dziennie i 250 słowach na minutę wychodzi \val{books.30.250} książek. To jest wynik idealnego roku, w którym każdy z 365 dni dostaje swoje trzydzieści minut. Takiego roku nie ma. Zwykle trzy dni z dziesięciu wypadają z różnych przyczyn.

Oto prosty sprawdzian. Trzydzieści minut w dwóch trzecich dni daje $30 \cdot 243 = 7\,290$ minut, a dwadzieścia minut każdego dnia to $20 \cdot 365 = 7\,300$ minut. Te dwa scenariusze są praktycznie równe. Dlatego realistyczny cel planujesz z wiersza \emph{dwadzieścia minut}: \val{books.20.250} książki, zaokrąglone w dół do 22, co jest celem z tabeli rocznej. Różnica między 22,8 a 22 to mniej niż jedna książka i stanowi bufor na sytuacje, których nie przewidziałeś.

Dla porównania, dwanaście książek rocznie wymaga około \val{need.12} minuty dziennie przy 250 słowach na minutę. To jest cel minimalny, który warto mieć na wypadek najgorszego roku. Cel dwudziestu czterech książek wymaga około \val{need.24} minut dziennie i jest celem na rok bez większych kryzysów.

\begin{zasada}
Cel roczny ustal jako wiersz tablicy odpowiadający minutom, które wytrzymasz w zły miesiąc, a nie w najlepszy. Plan oparty na najlepszym miesiącu łamie się przy pierwszym kryzysie, a plan oparty na złym miesiącu daje się utrzymać.
\end{zasada}

\section{Kiedy podnieść cel}

Cel wolno podnieść dopiero po trzech spokojnych miesiącach z rzędu. Spokojny miesiąc to taki, w którym zrealizowałeś co najmniej cztery z pięciu planowanych sesji tygodniowo i nie miałeś kryzysu z poprzedniej sekcji. Liczba cztery z pięciu jest tu wyborem, a nie wynikiem badań: jej zadaniem jest odróżnić miesiąc, w którym nawyk działa, od miesiąca, w którym przypadkiem wyszedł dobry tydzień.

Podnosisz cel o jeden wiersz tablicy, nie o więcej. Z dwudziestu minut na trzydzieści to skok o \val{books.30.250} minus \val{books.20.250}, czyli około 11,4 książki w teorii. To nie jest obietnica, tylko rachunek przy idealnym roku. Realny efekt będzie mniejszy, bo każdy nowy próg czasowy przynosi nowe wieczory, których dotąd nie planowałeś.

\begin{falle}[Pułapka]
Najczęściej podnosi się cel po dobrym okresie świątecznym albo po urlopie, czyli wtedy, gdy czasu było dużo, a nie wtedy, gdy nawyk był stabilny. Taki cel jest ustawiony na czas, którego zwykle nie ma. Jeśli podniesienie nastąpiło po dobrym tygodniu, zaczekaj do trzech spokojnych miesięcy, nawet jeśli wydaje ci się, że już możesz więcej.
\end{falle}

\section{Zamknięcie}

Ta książka zaczynała się od arytmetyki: ile książek mieści się w roku przy danej liczbie minut. Kończy się tym samym rachunkiem, ale z innej strony. Wiesz już, że dwadzieścia minut dziennie daje około dwudziestu dwóch książek przy typowej długości i szybkości, że porzucenie jest decyzją, a nie porażką, i że cel rośnie dopiero po stabilnym czasie.

Nie obiecuję, że po przeczytaniu tej książki będziesz czytać dwa razy więcej. Nie wiem tego i nie mogę tego sprawdzić z miejsca, w którym piszę. Mogę natomiast wskazać, co jest mierzalne: twój rejestr, twoje okna czasowe, twoje pięć pytań w grudniu. Jeśli po roku te dane wyglądają inaczej niż dzisiaj, książka zrobiła to, co mogła. Jeśli wyglądają tak samo, zmień plan, a nie siebie.

Zostaje jedna rzecz do zrobienia dziś, zanim zamkniesz tę stronę.

\begin{cwiczenia}
\exitem{Ile książek mieści się w roku przy 45 minutach dziennie i 200 słowach na minutę, jeśli typowa książka ma 80\,000 słów? Porównaj wynik z tablicą.}{Książka ma $80\,000 / 200 = 400$ minut czytania. Rok ma $45 \cdot 365 = 16\,425$ minut. Wynik to $16\,425 / 400 = 41{,}06$, czyli 41,1 książki. Zgadza się z wierszem tablicy dla 45 minut i 200 słów na minutę.}
\exitem{Ile minut dziennie potrzeba na dwanaście książek rocznie przy 250 słowach na minutę?}{Dwanaście książek to $12 \cdot 320 = 3\,840$ minut czytania rocznie. Podzielone przez 365 dni daje około 10,5 minuty dziennie, co jest wartością z tablicy.}
\exitem{Ustal cel miesięczny na miesiąc, w którym czytasz trudną książkę. Podaj liczbę książek i uzasadnij ją w jednym zdaniu.}{Wzorcowy wynik to jedna książka lub jedna trudna plus jedna lekka, a nie dwie trudne. Uzasadnienie: trudna książka ma wolniejszy postęp i zajmuje większość uwagi, więc drugi cel trudny zwiększyłby ryzyko porzucenia całego miesiąca.}
\exitem{Po miesiącu bez czytania wracasz. Co robisz pierwszego dnia i czego nie robisz?}{Pierwszego dnia czytasz minimum: pięć minut albo jedną stronę. Nie nadrabiasz zaległości i nie przenosisz celu na następny miesiąc. Zapisujesz w rejestrze przyczynę przerwy jednym słowem.}
\exitem{Dlaczego rejestr zlicza minuty i strony, a nie tylko skończone książki?}{Liczba skończonych książek ukrywa sesje, które nie zakończyły się książką, i zachęca do wybierania łatwych tekstów. Minuty i strony pokazują, czy nawyk działa, również wtedy, gdy żadna książka nie została zamknięta.}
\exitem{Po ilu spokojnych miesiącach z rzędu wolno podnieść cel i o ile wierszy tablicy?}{Po trzech spokojnych miesiącach z rzędu, o jeden wiersz tablicy, na przykład z dwudziestu minut na trzydzieści.}
\exitem{Wypisz pięć pytań grudniowej rewizji w kolejności, w jakiej je zadajesz.}{Zgodnie z sekcją: ile książek skończone i porzucone z powodem; które sesje były regularne; które kryzysy wystąpiły i który element planu zadziałał; czy książka zmieniła działanie i jaki był pierwszy krok; co zmienisz w przyszłym roku i dlaczego.}
\end{cwiczenia}

\zadanie{Ustal swój cel roczny z wiersza tablicy odpowiadającego minutom, które wytrzymasz w zły miesiąc, i zapisz go w rejestrze razem z dwoma pierwszymi wpisami z tabeli miesięcznej. Przez siedem dni prowadź rejestr według pól z Dodatku C. Miara sukcesu: po siedmiu dniach masz co najmniej pięć wpisów z pełnymi polami, a w każdym wpisie jest liczba minut, liczba stron i decyzja.}
